\chapter{无人机协同通信基础与问题框架}
\label{ch:foundations}

\section{引言}

无人机协同通信同时涉及物理连接、数据转发、语义表达和任务判断。不同研究层次使用的状态与评价指标有所区别，如果在建模时不作区分，容易把网络中传输的数据量、接收端获得的信息量和最终任务质量混为一谈。本章从无人机网络的基本运行过程出发，介绍后续研究所需的网络与能量模型、语义证据表示、在线优化和部分可观测决策方法，并建立三个研究问题之间的对应关系。

本章给出共同概念和分析工具，各研究章再按具体控制对象补充模型。路由研究以网络时隙为基本时间单位，关注多跳分组服务；语义选择以任务内的消息选择步骤为决策单位，关注不同证据的组合效果；交付控制则根据发送和反馈事件推进时间，关注候选消息的实际到达状态。三种时间组织方式分别适应各自的决策过程，在实验中不需要强行使用相同的采样频率。

\section{网络、资源与通信过程}
\subsection{动态连接与多跳数据回传}

将时刻 $t$ 的无人机通信网络记为有向图 $\mathcal G(t)=(\mathcal V,\mathcal E(t))$，其中 $\mathcal V$ 包含无人机节点和数据汇聚中心，$\mathcal E(t)$ 表示当前具备传输条件的链路。节点 $i$ 的位置和速度分别记为 $\bm x_i(t)$ 与 $\bm v_i(t)$，节点间距离为
\begin{equation}
d_{ij}(t)=\|\bm x_i(t)-\bm x_j(t)\|_2.
\label{eq:base:distance}
\end{equation}
链路是否可用既受到距离限制，也受到信道质量与当前资源分配的影响。邻居集合 $\mathcal N_i(t)=\{j:(i,j)\in\mathcal E(t)\}$ 因节点移动而改变，路由策略需要根据能够获得的邻居状态更新转发选择。

在简化无线模型中，链路速率可由带宽和接收信干噪比描述：
\begin{equation}
r_{ij}(t)=W_{ij}(t)\log_2\left(1+
\frac{P_i(t)h_{ij}(t)}{I_j(t)+N_j(t)}\right).
\label{eq:base:rate}
\end{equation}
其中 $W_{ij}$ 为可用带宽，$P_i$ 为发射功率，$h_{ij}$ 为信道增益，$I_j$ 和 $N_j$ 分别表示干扰与噪声功率。该模型说明了功率、带宽和传输能力之间的基本关系。提高功率能够改善特定链路的信干噪比，但对数形式意味着速率收益通常递减；同时，更高功率还会增加能量支出，并可能对其他接收端造成干扰。因此，局部功率选择需要放在整体资源约束下考虑。

多跳传输中，业务数据从源节点经过若干中继后到达汇聚中心。一次链路服务只使分组沿路径前进一跳，并不必然对应一份新的端到端业务数据完成交付。若统计所有链路上的转发量，较长路径可能带来更多累计传输字节；若评价唯一源数据的完成情况，则需要识别每个分组是否最终到达目的节点。这一区别对于解释吞吐量十分重要。第~\ref{rq1:chapter} 章采用累计转发量度量网络吞吐，并结合丢包率、能量分布和寿命评价路由性能。

\subsection{物理队列与能量状态}

数据队列记录已经到达但尚未获得服务的分组或比特。设节点 $i$ 在时隙 $t$ 开始时的队列长度为 $Q_i(t)$，新到达量为 $A_i(t)$，本时隙有效服务量为 $\mu_i(t)$，则基本更新关系为
\begin{equation}
Q_i(t+1)=\left[Q_i(t)-\mu_i(t)\right]^++A_i(t),
\qquad [x]^+=\max\{x,0\}.
\label{eq:base:queue}
\end{equation}
在多跳网络中，来自上游节点的转发还会进入下游队列，因此完整更新需要考虑入链路和出链路之间的关系。队列积压反映了通信需求尚未得到满足的程度，也为在线控制提供了可观测的拥塞信息。

能量状态反映节点后续能够执行的通信量。对一次持续 $\tau$ 的发送，基本发送能量为 $e^{\rm tx}=P^{\rm tx}\tau$。若节点初始能量为 $E_i(0)$，阶段 $t$ 的可计量支出为 $e_i(t)$，则剩余能量满足
\begin{equation}
E_i(t+1)=E_i(t)-e_i(t),\qquad 0\le e_i(t)\le E_i(t).
\label{eq:base:energy}
\end{equation}
不同研究需要明确能量核算的范围。无人机系统的总能耗通常还包括飞行、悬停、感知和计算等部分，发送能耗只是其中一项。本文根据各章研究对象限定核算边界：路由研究使用其网络仿真中的节点能耗模型，语义选择主要度量消息长度和传输代价，交付研究显式计算数据与确认的发送能耗。不同边界下的数值不能直接相加或横向比较。

在总能量相同的情况下，节点间支出分布也会影响网络持续性。若少数节点承担大部分转发任务，它们可能较早退出并破坏关键路径。反之，过度追求每个节点消耗相等，也可能让系统使用质量较差的链路或绕行过长路径。能量均衡需要与任务完成和传输效率共同评价，而不能独立理解为任何情况下都应最小化节点能耗方差。后续路由算法通过队列压力与剩余能量共同影响转发权重，交付算法则在可用语义候选范围内调整不同发送节点的使用程度。

\subsection{短通信窗口与反馈事件}

移动节点之间的有效接触往往具有时间限制。给定消息比特数 $B_m$ 与有效速率 $r_m$，其基本发送时长近似为 $B_m/r_m$。消息内容增加会延长发送时间，带宽不足或信道变差也会带来相同影响。若可用窗口长度为 $D_m$，则消息是否能够在窗口内完成成为通信决策的重要条件。语义选择可以利用预测窗口比较候选的风险，而交付控制还需要根据任务当前已经消耗的时间逐步缩小可行集合。

反馈使发送端能够更新对接收状态的判断，但反馈本身也是一个具有时长和不确定性的通信过程。一次数据发送至少涉及发送开始、发送结束、接收端生成确认、确认返回等事件。发送端在不同事件之后能够知道的信息不同：数据发送结束只说明本地尝试已经完成，收到累计确认才说明某个历史时刻接收端已获得哪些内容。区分这些时刻能够避免把本地动作与远端状态混淆。

\begin{figure}[htbp]
\centering
\begin{tikzpicture}[font=\small,>=Stealth]
\draw[->] (0,2.8) -- (12,2.8) node[right]{时间};
\draw[->] (0,0) -- (12,0);
\node[left] at (0,2.8){发送端};\node[left] at (0,0){接收端};
\fill (1,2.8) circle(1.5pt) node[above]{开始发送};
\fill (4,2.8) circle(1.5pt) node[above]{尝试结束};
\fill (4,0) circle(1.5pt) node[below,align=center]{载荷接收结果\\生成累计确认};
\draw[->,thick] (1,2.8)--node[above,sloped]{数据传输}(4,0);
\draw[->,thick,dashed] (4,0)--node[above,sloped]{反馈发送与延迟}(9.5,2.8);
\fill (9.5,2.8) circle(1.5pt) node[above]{收到确认};
\draw[decorate,decoration={brace,amplitude=4pt}] (4,3.6)--(9.5,3.6) node[midway,above=6pt,font=\footnotesize]{接收状态尚未完全确定};
\end{tikzpicture}

\bicaption{一次发送尝试中的数据事件与反馈事件}{Data and feedback events during one transmission attempt}
\label{fig:base:timeline}
\end{figure}

图~\ref{fig:base:timeline} 中，数据与反馈沿不同方向传递。在确认返回之前，发送端可能已经执行了下一次发送。因此，即使之后收到一个有效位图，该位图描述的也可能是较早时刻的接收集合。交付控制需要用位图的序号定位其对应状态，并把其后的发送记录重新纳入推断。反馈信息的时效性由其关联事件决定，而不只由消息到达的先后顺序决定。

\section{任务相关语义证据与质量度量}
\subsection{从局部观测到可通信证据}

设任务 $k$ 的局部观测为 $o_{i,k}$，任务查询或目标描述为 $q_k$。语义提取模块将二者映射为结构化证据 $e_{i,k}=\phi(o_{i,k},q_k)$。证据可以包含局部类别、置信度、观测质量、文本描述以及类别信念分布。与直接传输图像相比，这种表示把原始感知内容转换成任务可以利用的判断线索，使通信控制能够比较不同证据的成本和潜在作用。

本文采用五级证据可见性作为受控视觉任务的标签空间，分别对应清晰可见、大部分可见、不确定、大部分遮挡和完全遮挡。其目的在于分析局部视角是否包含支持特定视觉问题的证据，以及不同视角之间的信息交换如何改变中心判断。该任务由 VQAv2 与 COCO 图像及问题构造，但不直接评价标准视觉问答的自由文本答案生成。全图语义判断用作参考类别，局部节点只获得自身观测，二者的信息范围不同。

对固定类别空间 $\mathcal Y$，一个节点或中心的语义信念为
\begin{equation}
\bm b^{\rm sem}=(b_1,\ldots,b_{|\mathcal Y|}),\quad
b_y\ge0,\quad \sum_{y\in\mathcal Y}b_y=1.
\label{eq:base:semantic-belief}
\end{equation}
最大信念对应当前预测类别，熵 $H(\bm b^{\rm sem})=-\sum_y b_y\log b_y$ 描述分布的不确定程度。最大置信度和熵能够为候选选择提供信息，但它们本身不等于预测正确性。模型也可能对错误类别赋予较高置信度，因此算法需要通过训练与验证建立特征和任务质量之间的关系。

\subsection{语义粒度、冗余与互补}

同一份局部证据可以采用标签、判断陈述、摘要和完整证据等粒度表达。粒度较细时通常保留更多语义细节，同时增加消息长度。这里的“更多细节”并不保证“更大任务收益”：若短标签已经能够修正中心判断，继续发送同一观测的长摘要可能主要增加成本；若多个局部判断存在冲突，更详细的解释则可能帮助中心识别可靠证据。粒度选择需要结合当前接收端状态，而不宜固定采用最长或最短消息。

设已获得的证据集合为 $S$，任务质量函数为 $F_k(S)$，则新增证据 $m$ 的边际作用可写为
\begin{equation}
\Delta F_k(m\mid S)=F_k(S\cup\{m\})-F_k(S).
\label{eq:base:marginal}
\end{equation}
同一消息面对不同的已接收集合，其边际作用可能不同。一个与已有信息重复的证据可能几乎不改变判断；两个单独作用较弱的证据也可能共同消除歧义。因此，按固定分数独立排序与结合当前状态逐步选择，通常对应不同的控制能力。后续语义选择通过状态更新和策略网络处理这种顺序依赖，交付控制通过小规模接收集合的效用表处理候选之间的组合关系。

\begin{figure}[htbp]
\centering
\begin{tikzpicture}[
 x=1mm,y=1mm,>=Stealth,
 font=\fontsize{9bp}{11.5bp}\selectfont,
 box/.style={draw,line width=.45pt,align=center,text width=30mm,
   minimum height=13mm,inner sep=1mm},
 edge/.style={->,line width=.55pt},
 note/.style={align=center,text width=30mm,inner sep=0pt}]
\path[use as bounding box] (0,0) rectangle (148,-23);
\node[box] (a) at (16.25,-7.5) {局部图像\\与任务查询};
\node[box] (b) at (54.75,-7.5) {结构化证据\\与多粒度候选};
\node[box] (c) at (93.25,-7.5) {选择与交付\\资源、链路、反馈};
\node[box] (d) at (131.75,-7.5) {唯一证据融合\\与任务判断};
\draw[edge] (a.east) -- (b.west);
\draw[edge] (b.east) -- (c.west);
\draw[edge] (c.east) -- (d.west);
\node[note] at (54.75,-20) {内容是否有用};
\node[note] at (93.25,-20) {内容是否到达};
\node[note] at (131.75,-20) {判断是否正确};
\end{tikzpicture}

\bicaption{语义证据从局部形成到接收端使用的过程}{From local evidence extraction to receiver-side task inference}
\label{fig:base:semantic-pipeline}
\end{figure}

图~\ref{fig:base:semantic-pipeline} 区分了证据形成、候选筛选、物理交付和接收端融合。证据尚未被选择时，控制器需要判断它是否值得发送；证据已被选择但接收状态不明时，需要判断是否补发或等待；证据已经被确认接收时，重复成功发送不会产生新的内容。把这几种状态分开，能够更清楚地说明选择性通信与可靠交付之间的关系。

\subsection{分类质量与通信效率}

分类准确率统计最终预测与参考类别相同的任务比例。类别分布不均衡时，准确率可能掩盖少数类别的误判，因此语义选择实验还报告平衡准确率和宏平均 F1。有效信念质量定义为平均信念质量与平均送达概率的乘积，反映语义判断与链路送达条件的综合变化。它与只观察最终类别是否正确的准确率具有不同含义：信念分配改善或消息送达概率提高，都可能使该指标增加。

通信效率还需要考虑为形成上述结果付出的资源。消息 token 数描述语义表示长度，实际发送字节数还包括编码和协议头部；时延既可能来自数据发送，也可能来自等待反馈；能耗由功率和占用时长共同决定。对交付研究，单位正确任务通信能耗定义为全部通信支出与按时正确完成任务数的比值。这一指标同时反映正确交付和资源投入，但其分母变化也会影响结果，因此需要和正确完成率一起分析。

三个研究章根据各自控制对象定义质量指标。第~\ref{rq2:chapter} 章中的准确率评价固定语义轨迹上的最终类别判断，第~\ref{rq3:chapter} 章中的正确完成率还要求通信任务获得非空证据并在期限内完成。内容选择研究关注给定通信条件下的推理质量，可靠交付研究进一步纳入实际发送失败、反馈延迟和电量扣减。因此，两章的结果分别用于分析相应控制环节。

\section{在线优化与学习控制基础}
\subsection{李雅普诺夫漂移与虚拟队列}

在线网络控制需要在不知道未来到达、移动和信道状态的情况下选择当前动作。李雅普诺夫方法通过队列状态描述长期约束，把跨时隙问题转化为逐时隙可以处理的目标\cite{neely2010stochastic}。对长期平均能量约束 $\limsup_{T\rightarrow\infty}T^{-1}\sum_{t<T}\E[e_i(t)]\le\bar e_i$，可建立虚拟队列
\begin{equation}
Z_i(t+1)=\left[Z_i(t)+e_i(t)-\bar e_i\right]^+.
\label{eq:base:virtual-queue}
\end{equation}
与物理数据队列不同，$Z_i$ 不存放实际分组，而是记录能量支出超过参考额度的累积压力。它较大时表示该节点在过去一段时间中使用资源较多，应在后续决策中提高其能量成本。

由更新式的非负截断可以得到 $Z_i(t+1)\ge Z_i(t)+e_i(t)-\bar e_i$。对 $t=0,\ldots,T-1$ 求和可得
\begin{equation}
\frac{1}{T}\sum_{t=0}^{T-1}e_i(t)
\le \bar e_i+\frac{Z_i(T)-Z_i(0)}{T}.
\label{eq:base:finite-budget}
\end{equation}
该关系说明了虚拟队列与平均预算之间的联系：若队列增长相对于时间逐步消失，平均支出便接近约束范围。有限时域下，队列只是控制资源分配的状态量，仍需保留实际电量非负的硬约束。本文在第三项研究中借用虚拟队列记录跨任务压力，但不把有限任务回放直接解释为无限时域的稳定性证明。

取二次李雅普诺夫函数 $L(\bm Z)=\frac12\sum_iZ_i^2$，其条件期望漂移为 $\Delta(t)=\E[L(\bm Z(t+1))-L(\bm Z(t))\mid\bm Z(t)]$。在支出有界条件下，通过平方展开可以得到由常数项和 $\sum_iZ_i(t)(e_i(t)-\bar e_i)$ 构成的漂移上界。把当前服务收益的负值按系数 $V$ 加入，形成漂移加惩罚目标。队列项对应长期资源压力，收益项对应当前通信效果，两者之间的权重决定算法的取舍。第~\ref{rq1:chapter} 章结合具体无线速率与路由变量展开求解。

\subsection{动作约束与策略学习}

语义消息的任务作用较难直接写成简单解析函数，而候选特征、中心信念和预算状态又能够形成丰富的决策上下文，因此可采用强化学习构造选择策略。将任务状态记为 $s_h$，选择动作记为 $a_h$，即时奖励记为 $r_h$，策略目标为最大化有限步骤内的期望累计回报。状态更新反映证据被选择后的信念变化与预算扣减，奖励同时包含任务质量增量和通信代价。

实际候选集合会随步骤变化。已经选择的消息、超出剩余 token 预算的消息或不满足接收限制的消息不能再次执行。若网络输出候选分数 $z_a$，可通过可行集合 $\mathcal A(s)$ 构造掩码策略
\begin{equation}
\pi_\theta(a\mid s)=
\frac{\exp(z_a)\ind\{a\in\mathcal A(s)\}}
{\sum_{a'\in\mathcal A(s)}\exp(z_{a'})}.
\label{eq:base:masked-policy}
\end{equation}
停止动作通常保持可行，以保证其他候选都不可用时仍有明确的终止行为。动作掩码直接表达约束，能够减少策略把训练资源浪费在非法动作上；奖励成本则在可行候选之间调整偏好，二者作用不同。

近端策略优化通过限制新旧策略在已采样动作上的概率比，缓解一次更新过大带来的训练不稳定\cite{ppo}。其裁剪目标的基本形式为
\begin{equation}
L^{\rm clip}(\theta)=\E\left[
\min\left(\rho_h(\theta)\widehat A_h,
\operatorname{clip}(\rho_h(\theta),1-\epsilon,1+\epsilon)\widehat A_h\right)\right],
\label{eq:base:ppo}
\end{equation}
其中 $\rho_h$ 为新旧策略的动作概率比，$\widehat A_h$ 为优势估计，$\epsilon$ 为裁剪范围。价值网络提供后续收益估计，广义优势估计在偏差与方差之间取得折中\cite{gae}。这些训练机制属于通用工具，本文的具体方法设计主要体现在候选表示、链路与可靠性特征、状态更新、动作约束和任务奖励如何与通信问题结合。

\subsection{接收信念与部分可观测决策}

如果控制器不能直接观察真实接收集合，仅使用当前确认位图作为状态可能丢失重要信息。例如，过去多次发送虽然尚未得到确认，但已提高证据实际到达的概率。这种情况下，可以利用历史观测构造隐藏状态的条件分布，将过去的发送与反馈压缩为接收信念\cite{kaelbling1998pomdp}。设真实状态为 $x_h$、可观测历史为 $\mathcal H_h$，则
\begin{equation}
b^{\rm rec}_h(x)=\Prb\{x_h=x\mid\mathcal H_h\}.
\label{eq:base:reception-belief}
\end{equation}
这里的 $b^{\rm rec}$ 描述通信状态是否发生，而式~\eqref{eq:base:semantic-belief} 的 $\bm b^{\rm sem}$ 描述任务类别的判断。二者都是概率分布，但随机变量不同，不能互相替代。

给定动作引起的状态转移概率 $P(x'\mid x,a)$，在尚未获得新观测时，接收信念可以通过全概率公式预测：$\bar b(x')=\sum_xP(x'\mid x,a)b(x)$。新确认到达后，再依据该观测对候选状态的可能性进行修正。如果累计位图准确给出某一历史时刻的接收集合，该时刻的状态便退化为已知状态，而其后的未确认尝试仍需保留概率分支。第~\ref{rq3:chapter} 章利用候选数较小的特点，显式表示这些分支。

\begin{figure}[htbp]
\centering
\begin{tikzpicture}[font=\small,>=Stealth,
 box/.style={draw,rounded corners=2pt,align=left,text width=55mm,inner sep=8pt}]
\node[box] (a) at (0,0){\textbf{语义类别信念} $\bm b^{\rm sem}$\\随机变量：任务所属类别\\输入：已接收的证据内容\\输出：各类别的概率判断\\关注：这些证据意味着什么};
\node[box] (b) at (7.5,0){\textbf{接收状态信念} $b^{\rm rec}$\\随机变量：实际接收集合\\输入：发送历史与累计确认\\输出：各接收集合的概率\\关注：接收端可能拥有哪些证据};
\draw[->,thick] (b.south) -- ++(0,-.6) -| node[below,pos=.25,font=\footnotesize]{对不同接收集合的语义效用求期望} (a.south);
\end{tikzpicture}

\bicaption{语义类别信念与接收状态信念的区别}{Semantic class belief and reception-state belief}
\label{fig:base:beliefs}
\end{figure}

完整部分可观测规划需要考虑后续动作以及每种可能反馈，其分支数会随规划深度迅速增加。滚动前瞻通过限定预测范围，在当前状态下比较少量后续动作，只执行第一步，再利用实际返回的信息重新计算。这种方法保留了有限的未来影响，同时控制在线计算量。模型预测控制研究为受约束的滚动决策提供了基本分析视角\cite{mayne2000mpc}，但具体算法是否具备稳定性或最优性，还取决于终端条件、模型假设和求解方式，不能仅凭采用前瞻就直接获得相应保证。

\section{三项研究的统一问题框架}
\subsection{状态、动作与质量目标的对应}

从共同结构看，通信控制需要根据可观测状态，在资源约束允许的动作中选择能够改善当前或后续任务效果的动作。其一般形式可表示为
\begin{equation}
\max_\pi\ \E_\pi\left[\sum_h \left(R_h(s_h,a_h)-\lambda C_h(s_h,a_h)\right)\right],
\qquad a_h\in\mathcal A(s_h).
\label{eq:base:generic-control}
\end{equation}
式中 $R_h$ 表示研究层次对应的收益，$C_h$ 表示相应成本，$\mathcal A(s_h)$ 描述硬约束。该表达用于比较三个问题的结构，不意味着三个算法使用完全相同的状态转移、奖励函数或训练过程。表~\ref{tab:base:research-map} 给出了具体对应关系。

\begin{table}[htbp]
\centering\small
\bicaption{三个研究问题的控制对象与评价范围}{Control objects and evaluation scopes of the three research problems}
\label{tab:base:research-map}
\begin{tabularx}{\linewidth}{@{}p{19mm}YYY@{}}
\toprule
比较维度&能量感知路由&选择性语义通信&语义可靠交付\\
\midrule
决策对象&多跳数据分组&多来源、多粒度证据&给定候选的发送尝试\\
主要状态&邻居、队列、链路、能量&候选特征、中心信念、预算、历史&接收信念、累计确认、时间和电量\\
控制动作&下一跳、发送功率&消息选择、粒度、停止&新发、重传、等待、停止\\
主要约束&可达性、功率、队列与能量&token、步数、接收容量、送达门限&任务截止时间、节点电量、尝试次数\\
质量指标&吞吐、丢包、能量分布、寿命&分类质量、有效信念质量&按时正确完成率\\
资源指标&网络传输能耗与节点分布&token与通信代价&单位正确任务通信能耗\\
求解方式&队列驱动的局部在线优化&图结构策略与PPO&效用学习、概率滤波与有限前瞻\\
\bottomrule
\end{tabularx}
\end{table}

三项研究对长期和短期状态的使用也有所区别。路由中的数据队列与虚拟能量队列持续跨时隙更新；语义选择主要在单次任务内执行多步决策；交付控制既维护任务内的接收状态，也保留跨任务的节点电量和资源压力。后一种结构使当前任务中的低收益重传具有机会成本：它消耗的能量可能本可用于后续尚未开始的任务。

\subsection{研究衔接与实验边界}

第一项研究为后续提供的是对网络资源动态的认识和在线控制工具。它说明路由与功率决定局部传输能力，节点历史支出会改变未来通信机会。第二项研究在此基础上关注信息内容，把有限资源与证据的任务价值相联系。第三项研究进一步分析消息选择之后的实际交付，把随机发送结果和反馈信息纳入控制。三项研究沿资源分配、内容选择与实际交付逐步细化，共同描述通信机会如何转化为任务收益。

在方法接口上，路由或链路预测模块可以向语义层提供可用速率、传输时长与节点资源状态；语义选择模块可以为交付层提供较小规模的候选证据及其特征；交付控制则输出实际发送记录与确认状态，为后续资源更新提供依据。本文分别评价这些控制环节，以说明各项决策在相应条件下的作用。完整无人机系统还需要把不同时间尺度的状态同步、候选元数据交换和机载推理开销纳入统一运行过程，这些属于后续集成研究的内容。

实验中的训练、验证和测试承担不同作用。训练任务用于学习网络参数，验证任务用于选择控制权重或检查点，测试任务用于报告固定设置下的最终性能。不同网络随机性可以通过重复通信回放观察，但同一批视觉任务的多次回放不会增加独立图像样本数量。各研究章将明确其数据划分、运行次数和统计范围，防止把多次网络抽样误认为新增感知数据。

\section{本章小结}

本章介绍了网络连接、数据队列、能量状态、语义证据和反馈事件，说明了在线队列优化、掩码策略学习与接收信念控制的基本原理，并建立三项研究的共同框架。下一章从持续数据回传出发，研究能量感知的路由与功率联合决策。
